\scanpage{131}
est une alg. \deleted{à p.div.} \uline{graduée}, avec idéal $(S^{k'})^+$ \uline{hom.} à p.div. compatibles avec graduation, et $t^{k'}$ est un élém de degré $\delta$ \unsure{du dit}.

Alors les
\[
 A_n(\qquad)(k')=S^{k'}\{\{X_0,\ldots,X_n\}\}[dX_0,\ldots,dX_n]/(\textstyle\sum X_i-t,\sum dX_i)_{\mathrm{pd}}
\]
\inserted{une alg.} est \uline{bigraduée}, par le degré extérieur déjà envisagé ($2^{\text{ième}}$ degré) \inserted{(nul sur $S^{k'}$)}, et le degré total défini par
\[
\left\{\begin{array}{l}
\deg\operatorname{tot}X_r=\deg\operatorname{tot}dX_r=\delta\quad\deleted{\unsure{\deg}}\\
\deg\operatorname{tot}X_r^{(i)}=\unsure{\delta i|\delta};\quad\deg\operatorname{tot}_A\lambda=\deg\operatorname{tot}_{S^{k'}}\lambda\quad\big|\ \lambda\in S^{k'}
\end{array}\right.
\]
La diff. \deleted{\unsure{affecte}} \unsure{respecte} le degré total, $A_n(\quad)(k')$ est pour le degré tot une \uline{alg.} graduée sur $S^{k'}$ graduée commutative diff.; les puiss. divisées sont compatibles avec deux degrés.

Cette structure supplémentaire (deg tot.) est \deleted{\unsure{encore}} fonctorielle en $k'$, donc passe à $A^*(S,S^+,\gamma^*,t)$. Elle passe par suite aux $C_*^{\mathrm{DR}}(X;S,S^+,\gamma^*,t)$, et on retrouve la structure déjà connue de $C^*_{\mathrm{DR}}(X;S^{k'},(S^{k'})^+,\gamma_*^{k'},t^{k'})$ en prenant le $\operatorname{Hom}_k$ du \deleted{\unsure{$C$}} $C_*^{\mathrm{DR}}(X;S,S^+,\gamma^*,t)$ dans $k'$.

g) On prendra maintenant, \unsure{pour}
\[
\begin{array}{ll}
S=\deleted{\unsure{S}}S_k=k[u]&\text{donc }S^{k'}=k'\{\{T\}\}\\
\bar S=\bar S_k=k&\text{donc }\bar S^{k'}=k'
\end{array}
\]
\inserted{avec structure de coalg. \unsure{ordinaire}, \deleted{\illegible{1}} \unsure{sur $S_k$ la structure d'anneau}}
\[
\begin{array}{ll}
\bar S\longrightarrow S\quad\text{l'inclusion}&\text{donc }S^{k'}\longrightarrow\bar S^{k'}\quad\text{l'augm. canonique}\\
k\longrightarrow k[u]&k'\{\{T\}\}\longrightarrow k'
\end{array}
\]
\[
 \unsure{S^+}=k[u]/k\simeq k[u]^+
\]
\margin{130}
